\documentclass[../main]{subfiles}
\begin{document}
\chapter{Idea de proyecto}
\img{images/idea.jpg}{La idea de proyecto es aquello que nos acompaña desde el principio del proceso creativo. Sabemos el inicio del camino, esa idea, y sabemos el final del mismo: la forma en que nos gustaría materializar la idea }

\info{
Contenido}{
Idea de proyecto. Solución base optimizada.Identificación de necesidades. Identificación y organización de las diferentes soluciones.
}

\info{Solución base Optimizada}{
Una solución base optimizada para un proyecto es aquella que está diseñada específicamente para satisfacer los requisitos y objetivos del proyecto de la manera más eficiente posible. Esto implica que se ha realizado un análisis exhaustivo de los requerimientos y restricciones del proyecto, así como de las posibles soluciones tecnológicas disponibles.

Para optimizar una solución base, se deben considerar varios factores, como el rendimiento, la escalabilidad, la seguridad, la facilidad de mantenimiento, la interoperabilidad y el costo. La solución debe ser diseñada de manera que maximice el rendimiento del sistema y minimice el costo total de propiedad a largo plazo.

Además, una solución base optimizada debe ser escalable para permitir el crecimiento futuro del proyecto, ser segura para proteger los datos y los sistemas, y ser fácil de mantener y actualizar para minimizar el tiempo de inactividad del sistema.

En resumen, una solución base optimizada para un proyecto es una solución que satisface todos los requisitos y objetivos del proyecto de manera eficiente, económica y sostenible, y que es escalable, segura y fácil de mantener y actualizar.
}

\subsection{Idea de proyecto}
Una idea de proyecto es una propuesta inicial de una iniciativa que se tiene en mente para resolver un problema o aprovechar una oportunidad. Es una concepción inicial de lo que podría ser un proyecto, que puede ser desarrollada y refinada posteriormente.

Una idea de proyecto debe incluir los siguientes elementos básicos:
\begin{enumerate}
\item Descripción del problema u oportunidad a resolver: Es necesario identificar cuál es el problema u oportunidad a la que se quiere dar solución.

\item Objetivos y metas: Se deben establecer los objetivos y metas que se buscan alcanzar con el proyecto.

\item Alcance del proyecto: Es importante definir los límites del proyecto, es decir, qué se incluirá y qué no se incluirá dentro del proyecto.

\item Recursos necesarios: Se deben determinar los recursos que se requerirán para llevar a cabo el proyecto, como el presupuesto, personal, tiempo, materiales, entre otros.

\item Plan de acción: Se debe elaborar un plan de acción detallado, que incluya los pasos necesarios para llevar a cabo el proyecto y los plazos para cada uno de ellos.
\end{enumerate}
En resumen, una idea de proyecto es la primera etapa de cualquier iniciativa, es la concepción inicial de lo que podría ser un proyecto, donde se identifica el problema u oportunidad a resolver, se establecen los objetivos y metas, se define el alcance del proyecto.

\actividad{Actividad}
\begin{enumerate}
    \item Definir establecer grupos de trabajo para el proyecto.
    \item Definir idea de proyecto.
    \begin{enumerate}
        \item Realizar encuestas para encontrar problemáticas en la comunidad. Enumerar los problemas.
        \item Investigar por internet cuales son los principales problemas y necesidades de la sociedad.
        \item Encontrar diferentes soluciones a los problemas.
        \item Generar grupo de características o límites que deben cumplir las diferentes soluciones.
        \item Ponderar las diferentes opciones para elegir la mejor.
        \item Encontrar la solución base optimizada.
        \item Encontrar 4 alternativas más al proyecto.
    \end{enumerate}
    \item Definir la idea de proyecto para su proyecto utilizando el listado explicado en la teoría.
\end{enumerate}





\end{document}